% GENERATED FILE. Edit canonical lessons, then run npm run generate:books.
% canonical-source-hash: fnv1a64:ddd9017247023f42
% canonical-lessons: KA-S152-full-stop-and-comma, KA-S153-question-mark, KA-S154-other-marks, KA-R77-marks-recall, KA-W77-gotta-controlled-question, KA-W77-a1-timed-production

\chapter{Marks on the Page, Your First Question, and Timed Writing}
\label{ch:ka-viraama-chihne}
\hlchaptermodality{77}

\begin{chapteropening}
\textbf{By the end of this chapter:} \emph{I can name Kannada punctuation, compose a familiar question, and complete the A1 writing paper within 20 minutes.}

Complete the chapter's final lesson: fill a six-field form and write a 20–30 word named-reader message within the A1 paper's 20-minute limit.
\end{chapteropening}

\section[. , (full stop, comma)]{. and , — two marks you have been reading all along}

\label{lesson:KA-S152-full-stop-and-comma}



\begin{quote}
Before the new ones: \kn{ಟ} — what sound does it carry?

Now look at the end of any Kannada line in the last chapter. Something is sitting there that no lesson has ever named.
\end{quote}

\subsection*{Script you'll notice: the full stop and the comma}
\textbf{.} — the full stop, which ends a sentence.

\textbf{,} — the comma, which separates items inside one.

\textbf{Kannada borrowed both of these whole from the Latin alphabet}, shape and job together, and uses them the way English does. They sit on the baseline, below the letters, rather than hanging from anything.

You have been reading them since the first page:

\begin{quote}
\kn{ಇದು} \kn{ಬಾಗಿಲು}. \kn{ಕುರ್ಚಿ} \kn{ಇಲ್ಲಿ} \kn{ಇದೆ}.
\end{quote}

Two sentences, two full stops, and until now nothing in this book had said what they were.

\textbf{Older Kannada printing used an upright stroke to end a sentence instead.} You are unlikely to write one, and you may meet one in an old book or on a sign that was set once and never reset. Knowing it is a sentence-ender is the whole of what it asks of you.

\subsection*{Writing: . and ,}
Copy the two marks at the foot of a line of Kannada, once each. Put the full stop tight against the last letter, with the space after it rather than before.

\begin{quote}
This book does not yet tell you \textbf{where to start the character or which way to travel}. That is a real thing, taught with real variation from school to school, and it is not written down here until it can be written down with a source. Copying what is in front of you needs no such source, and it is how the shape gets into your hand in the meantime.
\end{quote}

\subsection*{Your turn}
\begin{itemize}
\raggedright
  \item \emph{Look:} at the two lines above and find the mark that ends each one
\end{itemize}

\begin{quote}
\kn{ಇದು} \kn{ಬಾಗಿಲು}.  ·  \kn{ಕುರ್ಚಿ} \kn{ಇಲ್ಲಿ} \kn{ಇದೆ}.
\end{quote}

\begin{itemize}
\raggedright
  \item \emph{Say it:} which of the two marks ends a sentence and which divides one
  \item \emph{Look:} back at any page of this book and find a full stop on it
\end{itemize}

\begin{tcolorbox}[breakable,colback=teal!4,colframe=teal!35!black,title={Before you move on}]
What does a full stop do? (Ends a sentence.) What does a comma do? (Separates items inside one.) Where did Kannada get both of them? (From the Latin alphabet, shape and job together.)
\end{tcolorbox}
\section[? (question mark)]{? — one mark, at the end}

\label{lesson:KA-S153-question-mark}



\begin{quote}
Before the new one: what does a full stop do, and where does it sit?

Today's mark does the same job for a different kind of sentence.
\end{quote}

\subsection*{Script you'll notice: the question mark}
\textbf{?} — the question mark.

Like the full stop, Kannada takes it from the Latin alphabet whole. It ends a question the way a full stop ends a statement, and it goes in \textbf{one place only: the end}.

\noindent
\begin{tabularx}{\linewidth}{@{}>{\raggedright\arraybackslash}X>{\raggedright\arraybackslash}X@{}}
\toprule
\textbf{sentence} & \textbf{mark} \\
\midrule
a statement & . \\
a question & ? \\
\bottomrule
\end{tabularx}

\textbf{Nothing marks the opening of a question.} Some languages put a second, upside-down mark at the front of one; Kannada does not, and neither does English. A reader meets the question mark at the end or not at all, which means the words themselves have to carry the question until then.

That is worth noticing, because it changes how you read aloud: the rise in your voice has to come from recognising a question word, not from a mark you have not reached yet.

\subsection*{Writing: ?}
Copy the mark once at the end of a line, tight against the last letter. Then write a full stop beside it and look at the two.

\begin{quote}
This book does not yet tell you \textbf{where to start the character or which way to travel}. That is a real thing, taught with real variation from school to school, and it is not written down here until it can be written down with a source. Copying what is in front of you needs no such source, and it is how the shape gets into your hand in the meantime.
\end{quote}

\subsection*{Your turn}
Say these aloud:

\begin{itemize}
\raggedright
  \item which mark ends a question and which ends a statement
  \item how many marks a Kannada question carries (one)
  \item where in the sentence that one mark goes
\end{itemize}

\begin{tcolorbox}[breakable,colback=teal!4,colframe=teal!35!black,title={Before you move on}]
Where does a Kannada question mark go? (At the end.) Is there a second one at the front? (No.) What has to tell you a question is coming before you reach it? (The words.)
\end{tcolorbox}
\section[: ( ) (colon, brackets, quotes, dash)]{: ( ) — the rest of the set, for reading}

\label{lesson:KA-S154-other-marks}



\begin{quote}
Before the new ones: where does a question mark go, and how many does a sentence get?

Four more marks, and this time the job is recognising them rather than reaching for them.
\end{quote}

\subsection*{Script you'll notice: the colon, the brackets, the quotes and the dash}
\noindent
\begin{tabularx}{\linewidth}{@{}>{\raggedright\arraybackslash}X>{\raggedright\arraybackslash}X@{}}
\toprule
\textbf{mark} & \textbf{what it tells a reader} \\
\midrule
: & a list or an explanation follows \\
( ) & an aside you could lift out \\
" " & somebody's exact words \\
— & a break in the sentence \\
\bottomrule
\end{tabularx}

\textbf{Kannada borrows the whole Latin set}, which is the useful fact: there is no second system to learn and no Kannada-shaped equivalent to look for. A mark you recognise from English does the job you expect it to do.

\textbf{These are for reading.} A form, a notice or a printed page will put them in front of you, and at this level the demand is knowing what each one signals rather than deciding where to place one yourself. Placing them is a writing habit that comes later, and it comes from reading a great many of them first.

\subsection*{Writing: : and ( )}
Copy the colon and the pair of brackets once each. The brackets come in twos — one opens, one closes, and a line with only one of them is a line with something missing.

\begin{quote}
This book does not yet tell you \textbf{where to start the character or which way to travel}. That is a real thing, taught with real variation from school to school, and it is not written down here until it can be written down with a source. Copying what is in front of you needs no such source, and it is how the shape gets into your hand in the meantime.
\end{quote}

\subsection*{Your turn}
\begin{itemize}
\raggedright
  \item \emph{Say it:} which mark tells you a list is coming
  \item \emph{Say it:} which pair you could lift out of a sentence without breaking it
  \item \emph{Say it:} which mark means the words are somebody else's
  \item \emph{Look:} at any printed page and find two of the four
\end{itemize}

\begin{tcolorbox}[breakable,colback=teal!4,colframe=teal!35!black,title={Before you move on}]
Which of these four arrive in pairs? (The brackets, and the quotation marks.) Where did Kannada get the set? (From the Latin alphabet, whole.) Are you being asked to place them or to recognise them? (To recognise them.)
\end{tcolorbox}
\section[(marks)]{(marks) — which ones you write, and which ones you only read}

\label{lesson:KA-R77-marks-recall}



\begin{quote}
Close the three lessons before this one. Name the mark that ends a Kannada statement, with nothing in front of you.
\end{quote}

\subsection*{Script: the whole set, sorted by what it asks of you}
\noindent
\begin{tabularx}{\linewidth}{@{}>{\raggedright\arraybackslash}X>{\raggedright\arraybackslash}X@{}}
\toprule
\textbf{you write these} & \textbf{you read these} \\
\midrule
. , ? & : ( ) " " — \\
\bottomrule
\end{tabularx}

\textbf{Every one of them came from the Latin alphabet}, so there was never a second system to learn — which is why this chapter is short and why nothing in it looked unfamiliar.

The split in that table is the part worth keeping. A reader at this level puts full stops, commas and question marks on the page themselves, and recognises the rest when a notice or a form puts them there.

\subsection*{Your turn}
\begin{itemize}
\raggedright
  \item \emph{Say it:} the three marks you write yourself
  \item \emph{Say it:} where a question mark goes, and how many a question gets
  \item \emph{Say it:} what an older Kannada book might use at the end of a sentence
  \item \emph{Look:} at the reading passage in the chapter before this one and name every mark in it
\end{itemize}

\begin{tcolorbox}[breakable,colback=teal!4,colframe=teal!35!black,title={Before you move on}]
Which three marks do you write? (\textbf{.} \textbf{,} \textbf{?}) Which mark opens a question? (None.) Where did the whole set come from? (The Latin alphabet.)
\end{tcolorbox}
\section[\texorpdfstring{\kn{ನಿಮಗೆ} \kn{ಗೊತ್ತಾ}? (nimage gottā?)}{ನಿಮಗೆ ಗೊತ್ತಾ? (nimage gottā?)}]{Your first Kannada question from meaning}

\label{lesson:KA-W77-gotta-controlled-question}



\begin{quote}
Close the earlier lessons. Recall how Kannada marks the person who knows, the word meaning \textbf{known}, the ending that turns a statement into a yes-or-no question, and the mark that closes a question. Do not write yet.
\end{quote}

\subsection*{Writing — assemble one familiar question}
From meaning alone, write \textbf{Do you know?} There is no word bank, romanization, or copyable sentence. Choose the familiar pieces, put them in Kannada order, change the know-word into a yes-or-no question, and add the question mark. This first composition is untimed: finish before opening the answer key.

\subsection*{Your turn — repair one difference}
Check the knower, know-word, long question ending, word boundary, and punctuation separately. Repair the first difference and write the whole question once more from meaning.

\begin{tcolorbox}[breakable,colback=teal!4,colframe=teal!35!black,title={Before you move on}]
What made this composition rather than copying? (\textbf{You selected, changed, and ordered known pieces from meaning, with no visible Kannada answer.})

One controlled line proves this rung, not the later timed A1 writing task.
\end{tcolorbox}
\section[(timed A1 writing)]{Your first complete Kannada A1 writing paper}

\label{lesson:KA-W77-a1-timed-production}



\begin{quote}
This is an attempt, not a model. You will write the answers, and this page will not show you any copyable Kannada answer.

Put away romanization, dictionaries, translators, and spell-checkers. Set one timer for \textbf{20 minutes}. Keep it running through both tasks. Write in Kannada script. Ordinary legible handwriting is enough; decorative calligraphy is not scored.
\end{quote}

\subsection*{Writing — timed assessment production}
\#\#\# Task 1 — practical form: 7 minutes, 40 points

Complete \textbf{all six fields in Kannada script} for a neighbourhood language exchange. Short answers are enough.

\begin{enumerate}
\raggedright
  \item Your name
  \item A language you speak
  \item One greeting you would use
  \item One drink you would ask for
  \item One place-question word
  \item One polite closing you would use
\end{enumerate}

At seven minutes, move on even if a field is unfinished. Do not restart the timer and do not polish the form with time meant for the message.

\#\#\# Task 2 — message: 13 minutes, 60 points

\textbf{Reader:} Mira. \textbf{Purpose:} help Mira recognise you and arrange a first conversation tomorrow.

Write \textbf{20–30 words in Kannada script}. Greet Mira, give your name, say that you speak Kannada, mention the drink you would ask for, ask where Mira is, ask one familiar yes-or-no question, and close politely. The instructions give content, not sentences: choose and connect the Kannada yourself.

When the 20-minute timer rings, \textbf{stop}, even if a line is unfinished. Record the number of completed form fields and the message word count. An honest timed attempt is more useful than a perfect untimed one.

\begin{tcolorbox}[breakable,colback=teal!4,colframe=teal!35!black,title={Before you move on}]
Only after you have stopped, review the same attempt. Do not rewrite it yet. Mark one strength and one next step under each assessment criterion:

\begin{enumerate}
\raggedright
  \item \textbf{Task fulfilment and relevant content} — six fields are attempted; the named reader, conversation purpose, and requested message points are present.
  \item \textbf{Comprehensibility and organisation} — Mira can follow the message and see how its ideas connect.
  \item \textbf{Kannada vocabulary, grammar, and register} — familiar words and sentence patterns carry the intended meaning in a suitable everyday tone.
  \item \textbf{Kannada orthographic control} — vowel-sign placement, consonant clusters, word spacing, and punctuation remain readable.
\end{enumerate}

Keep the paper. On another day, repeat the same two-task shape with new personal details. The goal is not to memorise this prompt; it is to make a complete Kannada attempt calmly inside the real limit.
\end{tcolorbox}
